% !TeX program = pdflatex
% papers/inteiros.tex — O DEGRAU N → Z: reversibilidade da soma, palavra em 8 bits.
%
% Auto-contido: cd papers && pdflatex inteiros.tex
\documentclass[11pt,a4paper]{article}
\usepackage[utf8]{inputenc}
\usepackage[T1]{fontenc}
\usepackage[portuguese]{babel}
\usepackage{amsmath,amssymb,amsthm}
\usepackage[margin=2.4cm]{geometry}
\usepackage{booktabs}
\usepackage{xcolor}
\usepackage[hidelinks]{hyperref}

\newtheorem{teorema}{Teorema}
\newtheorem{proposicao}[teorema]{Proposição}
\newtheorem{corolario}[teorema]{Corolário}
\newtheorem{definicao}{Definição}
\newtheorem{observacao}{Observação}

\newcommand{\code}[1]{\texttt{\small #1}}
\newcommand{\FF}{\mathbb{F}}
\newcommand{\abs}[1]{\lvert #1\rvert}
\newcommand{\dg}{^{\dagger}}
\newcommand{\Word}{\mathcal{W}}
\newcommand{\sepcol}{\quad\penalty0$\big|$\quad\penalty0}
\newcommand{\colunas}[3]{%
  \par\smallskip\noindent{\footnotesize\raggedright
  \textbf{prova} (Algébrico): #1\sepcol
  \textbf{realiza} (aqui): #2\sepcol
  \textbf{verifica}: #3\par}\smallskip}

\definecolor{cinza}{gray}{0.35}
\newcommand{\medido}[1]{\par\medskip\noindent{\small\color{cinza}\textbf{Medido.} #1}\par\medskip}

% ─────────────────────────────────────────────────────────────────────────────
%  papers/gkcapa.tex — a capa do Reino, mínima, para os papers em `article`.
%
%  O estilo.tex completo é do `report` (títulos de capítulo, skak, …). Aqui fica
%  só o que a capa precisa, para o paper continuar a compilar sozinho com
%  pdflatex. O tradutor próprio (tests/tex) carrega o estilo.tex da raiz / o
%  symlink papers/estilo.tex e avalia o mesmo `\gkcapa`.
% ─────────────────────────────────────────────────────────────────────────────
\definecolor{ouro}{HTML}{B8912F}
\definecolor{ouroclaro}{HTML}{F3E9CE}
\definecolor{tinta}{HTML}{1E1B16}
\definecolor{regua}{HTML}{6E6858}
\providecommand{\gknota}{\fontsize{7.62}{11.05}\selectfont}
\providecommand{\gktexto}{\fontsize{10.50}{15.22}\selectfont}
\providecommand{\gksub}{\fontsize{12.33}{17.87}\selectfont}
\providecommand{\gksec}{\fontsize{14.47}{20.98}\selectfont}
\providecommand{\gkcap}{\fontsize{16.99}{24.63}\selectfont}
\providecommand{\gktit}{\fontsize{23.42}{33.95}\selectfont}
\providecommand{\Prj}{\mathbb{P}}
\providecommand{\Trf}{\mathcal{T}}
\providecommand{\gkcapa}[3]{%
\title{\vspace{-0.6cm}
{\color{ouro}\rule{\textwidth}{1.4pt}}\\[6mm]
\textbf{\gktit\color{tinta} Reino Dourado}\\[3mm]{\gksec\itshape\color{regua} Golden Kingdom}\\[8mm]
{\gkcap\scshape\color{ouro} #1}\\[5mm]
{\gksub\color{regua} #2}\\[2mm]
{\gktexto\color{regua}\itshape #3}\\[7mm]
{\color{ouro}\rule{\textwidth}{0.6pt}}\\[9mm]{\gknota\color{regua}\begin{minipage}{\textwidth}\setlength{\parindent}{0pt}%
\textbf{Aviso de Isenção Algorítmica:} O universo, as leis e as entidades contidas nestas páginas ---
incluindo felinos matriciais, aves de rapina algébricas e amuletos de controle de fluxo --- são
construções de pura ficção formal. Esta obra não descreve a realidade física, não serve como
engenharia reversa do cosmos e não constitui doutrina religiosa. Qualquer semelhança com bugs reais do seu
sistema operacional é mera coincidência (ou falta de reza). Prossiga por sua própria conta e
risco.\end{minipage}}}
\author{}\date{}}


\begin{document}
\gkcapa{Os inteiros a partir dos naturais}
       {$\mathbb{Z}=\mathbb{N}^{2}/{\sim}$ --- classe de equivalência,
        soma cruzada $a+d=b+c$}
       {oito leis, pares, soma cruzada, $\nu$, memória por duplicidade}
\maketitle

%======================================================================
\begin{abstract}
Degrau $k{=}0$ da torre (\code{analitico thm:tecidos}): construção dual
$T+T^{*}$ --- $\mathbb{N}^{2}/{\sim}$ com $(a,b)\sim(c,d)\iff a+d=b+c$.
O inteiro é a \textbf{classe} $[(a,b)]$; $\Phi([(a,b)])=a-b$.
\end{abstract}

%======================================================================
\section*{Onde este documento vive}

\noindent\textbf{Camada N--R}, não a Parte Álgebra, não \(\operatorname{Star}(\mathcal{U})\).
\[
\boxed{\;
\operatorname{Star}(\mathcal{U})=\mathcal{D}
\quad\neq\quad
\operatorname{Star}(\mathbb{Z})=\varnothing
\quad\neq\quad
\text{o que a Física constrói}
\;}
\]
Passagem por \textbf{opostos}, não por \(\mathcal{D}\). \(\nu(z)=-z\neq\mathcal{D}\).

Degrau $k=0$ da torre $T_{k+1}=T_k+T_k^{*}$
(\code{analitico thm:tecidos}) --- o \textbf{segundo} da construção, que se lê
num sentido só:

\[
\boxed{\;
\underbrace{\FF_2}_{\text{um bit}}
\;\xrightarrow[\text{3 dobras duais}]{\ \code{naturais}\ }\;
\underbrace{F_8\equiv\Word_8}_{\{0,\ldots,255\}}
\;\xrightarrow[\text{TRANSPORTE}]{}\;
\mathbb{N}
\;\xrightarrow[\text{soma}]{\ \code{inteiros}\ }\;
\mathbb{Z}
\;\xrightarrow[\text{produto}]{\ \code{racionais}\ }\;
\mathbb{Q}
\;\xrightarrow[\text{corte}]{\ \code{reais}\ }\;
\mathbb{R}\;}
\]
\noindent{\footnotesize Lê-se da esquerda para a direita: é o sentido em que se
\textbf{constrói}. Cada paper recebe o objecto do degrau à sua esquerda e
entrega o do degrau à sua direita; nenhum redefine o que recebeu.\par}

\noindent A cadeia completa, comum aos três papers
--- a mesma nos três:

\begin{center}\small
\begin{tabular}{@{}clllll@{}}
\toprule
$k$ & degrau & ganha volta & paga & dual do andar (e o seu Fix) & paper \\
\midrule
$-1$ & $F_1\!\to\!F_2\!\to\!F_4\!\to\!F_8$ & a \textbf{largura} (3 dobras)
     & a \textbf{ordem} (car.\ $2$) & $x^{*}$ (2.ª raiz); Fix $=$ andar de baixo
     & \code{naturais.tex} \\
$0$  & $\mathbb{N}\!\to\!\mathbb{Z}$ & a \textbf{soma}
     & o sinal e a boa ordem & $\iota(a,b)=(b,a)\Rightarrow$ oposto; Fix $=\{0\}$
     & \code{inteiros.tex} \\
$1$  & $\mathbb{Z}\!\to\!\mathbb{Q}$ & a \textbf{multiplicação}
     & $0^{-1}$ não existe & $\iota(a,b)=(b,a)\Rightarrow$ inverso; Fix $=\{\pm1\}$
     & \code{racionais.tex} \\
$2$  & $\mathbb{Q}\!\to\!\mathbb{R}$ & o \textbf{corte} (completação)
     & a \textbf{enumerabilidade} & $(A,B)\mapsto(B,A)$; Fix: o corte que fecha
     & \code{reais.tex} \\
\bottomrule
\end{tabular}
\end{center}
\noindent{\footnotesize A numeração é a de \code{arquitetura sec:descida},
onde $k=0$ é $\mathbb{N}\!\to\!\mathbb{Z}$; a torre binária é \emph{anterior} a
esse degrau, e por isso $k=-1$. O passo é sempre $T_{k+1}=T_k+T_k^{*}$
(\code{analitico thm:tecidos}). \textbf{Entre $k=-1$ e $k=0$ não há dobra:
há o TRANSPORTE} --- $a+b=(a\oplus b)+2(a\wedge b)$
(\code{naturais.tex thm:transporte}) ---, e é ele que leva o suporte $\Word_8$
do topo da torre binária às operações de $\mathbb{N}$. Nos degraus $k\ge0$
duplica-se e \textbf{quocienta-se}; em $k=-1$ duplica-se e \textbf{não} se
quocienta (a segunda cópia é multiplicada por $\sigma$).\par}

\noindent E o dual, que em cada degrau tem \emph{duas} partes:

\[
\boxed{
\begin{array}{lll}
 & \textbf{dual do CENTRO (Lei 1)} & \textbf{dual do ANDAR (o que o degrau acrescenta)}\\[3pt]
k=-1 & \nu(x)=-x=x\ \text{(degenera)} & x^{*}=(a_0{+}a_1)+a_1\sigma\ \Rightarrow\ T=x\oplus x^{*},\ N=x\otimes x^{*}\\[2pt]
k=0  & \nu(z)=-z & \iota(a,b)=(b,a)\ \Rightarrow\ \text{oposto aditivo}\\[2pt]
k=1  & \nu(a,b)=(-a,b) & \iota(a,b)=(b,a)\ \Rightarrow\ \text{inverso multiplicativo}
\end{array}}
\]
\noindent{\footnotesize \textbf{O mesmo \emph{swap} induz o oposto num degrau e
o inverso no seguinte} --- é o dado que muda de papel, não a fórmula. E em
$k=-1$ o dual do centro degenera ($-x=x$), pelo que o dual tem de vir de outro
sítio: da segunda raiz de $\sigma^2+\sigma+\lambda$. Nenhuma destas involuções
é o dual linear $V\!\cong\!V^{**}$ da Lei~2, nem o dual quadrático
$s\dg=-1/s$ do \emph{Corpo Algébrico}: são quatro noções, e não se confundem.\par}

\noindent\textbf{O que RECEBE do degrau anterior.} $\mathbb{N}$ entra aqui
por Peano (Def.~\ref{def:peano}) --- \emph{não} se constrói neste paper. O
degrau anterior (\code{naturais.tex}, $k=-1$) ergue a torre binária
$F_1\!\to\!F_8$ e entrega o \textbf{suporte}: o envelope $\Word_8$ de
§\ref{sec:palavra8} é o suporte de $F_8$, com os andares de baixo como
prefixos. A passagem desse suporte para as operações de $\mathbb{N}$ é o
\textbf{transporte}, e não uma dobra --- por isso $\mathbb{N}$ chega aqui como
hipótese e não como construção.

\medskip\noindent\textbf{O que ENTREGA ao degrau seguinte.} Ao
\code{racionais.tex} ($k=1$) entrega $\mathbb{Z}$ e o \emph{mesmo} molde de
pares: lá o par volta, com produto cruzado em vez de soma cruzada, e o mesmo
\emph{swap} $\iota(a,b)=(b,a)$ passa a induzir o inverso em vez do oposto.
Entrega também a descida $\Word_8\to E_{16}$ (§\ref{sec:palavra8}).
\textbf{Molde dual} (Lei~4, $T+T^{*}$): duplicar, ler nos dois lados,
quocientar. Em $\mathbb{N}$ a soma é total e a subtracção é \textbf{parcial}
(\code{teoria thm:fibra-por-andar}). Em $\mathbb{Z}$ a subtracção torna-se
\textbf{total}; paga-se o \textbf{sinal} e a \textbf{boa ordem} de $\mathbb{N}$
--- \emph{não} a ordem total de $\mathbb{Z}$, que permanece via
$\Phi$ após o quociente (Def.~\ref{def:espaco}). A folha $\iota(a,b)=(b,a)$ realiza o dual no
espaço; $\nu(z)=-z$ no quociente (\code{algebrico cor:cadeia},
$\mathbb{Z}=\mathbb{N}+{}$sinal).

\medskip\noindent\textbf{Dois circuitos em paralelo} (padrão de toda a série):
\[
\mathbb{N}\xrightarrow{T+T^{*}}\mathbb{N}^{2}
\xrightarrow[\mathrm{Cruz}]{\mathrm{Dir}}\xrightarrow{\,q\,}\mathbb{Z}
\qquad
I\xrightarrow{p}\mathbb{Z}^{2}\xrightarrow{|p^{-1}|}G_{\mathrm{geom}}.
\]
O \emph{Racionais} (\code{racionais sec:torre}) trata o degrau seguinte com o
mesmo molde e produto cruzado.

\medskip\noindent\textbf{O que este documento RECEBE e o que ACRESCENTA}
--- a mesma estrutura que no \emph{Racionais} e no \emph{Corpo Algébrico}
(\code{algebrico def:espaco}):

\[
\boxed{
\begin{array}{lll}
\textbf{RECEBE do centro} & \textbf{onde está lá} & \textbf{o que se faz aqui}\\[3pt]
\text{a torre }T_{k+1}=T_k+T_k^{*} & \code{analitico thm:tecidos} & \text{instancia o degrau }k{=}0\\[2pt]
\text{as oito leis} & \code{analitico sub:oitoleis} & \text{Leis 0--4,7 explícitas; 5--6 recebidas}\\[2pt]
\text{o par dual e Dir/Cruz} & \code{algebrico def:espaco} & \text{soma cruzada em }\mathbb{N}^{2}\\[2pt]
\text{Peano e cúspide} & \code{algebrico thm:cuspide} & \text{prova Peano; cúspide realiza}\\[6pt]
\textbf{ACRESCENTA} & \multicolumn{2}{l}{\text{---e é só isto---}}\\[3pt]
\text{a REPRESENTAÇÃO por pares} & \text{Def.~\ref{def:espaco}} & V_{\mathrm{int}}/{\sim}=\mathbb{Z}\\[2pt]
\text{a EQUIVALÊNCIA }{\sim} & \text{Prop.~\ref{prop:equiv}} & a+d=b+c\ \text{antes do quociente}\\[2pt]
\text{o ISOMORFISMO }\Phi & \text{Thm.~\ref{thm:z}} & \Phi([(a,b)])=a-b\\[2pt]
\text{a FOLHA e }\nu & \text{§\ref{sec:sinal}} & \text{inversão }\nu(z)=-z;\ \Phi(\nu(x))=-\Phi(x)\\[2pt]
\text{dois circuitos }q\parallel p & \text{§\ref{sec:memoria}} & G_{\mathrm{alg}}\ \text{vs.}\ G_{\mathrm{geom}};\ \Word_8
\end{array}}
\]

%======================================================================
\section{As oito leis --- marcos da construção}\label{sec:leis}

\noindent Onde cada lei vive na cadeia --- tabela comum aos três papers
--- a mesma nos três:

\begin{center}\footnotesize
\begin{tabular}{@{}cp{2.75cm}p{2.75cm}p{2.75cm}p{2.75cm}@{}}
\toprule
lei & $k{=}-1$: $F_1\!\to\!F_8$ & $k{=}0$: $\mathbb{N}\!\to\!\mathbb{Z}$ & $k{=}1$: $\mathbb{Z}\!\to\!\mathbb{Q}$ & $k{=}2$: $\mathbb{Q}\!\to\!\mathbb{R}$ \\
\midrule
\textbf{0} divisão do zero
 & $x\oplus x=0$; só o $0$ sem inversa \hfill\textsc{expl.}
 & subtracção parcial; folha $[(1,0)]\oplus[(0,1)]=0$ \hfill\textsc{expl.}
 & $0^{-1}$ não existe; $0\!\leftrightarrow\!\infty$ \hfill\textsc{mista}
 & o \textbf{buraco}: corte sem racional em cima \hfill\textsc{expl.} \\
\textbf{1} dual ($1\dg=-1$)
 & \emph{degenera}: $-x=x$ \hfill\textsc{expl.}
 & $\iota$ no par; $\nu(z)=-z$ \hfill\textsc{expl.}
 & $\nu(a,b)=(-a,b)$ \hfill\textsc{expl.}
 & $\nu(x)=-x$; do andar: $(A,B)\!\mapsto\!(B,A)$ \hfill\textsc{expl.} \\
\textbf{2} bidual
 & $x^{**}=x$; automorfismo \hfill\textsc{expl.}
 & Dir/Cruz; $R^{4}=\mathrm{id}$ \hfill\textsc{mista}
 & $V\cong V^{**}$ canónico \hfill\textsc{canón.}
 & $B=\mathbb{Q}\setminus A$: complementar $2\times$ \hfill\textsc{expl.} \\
\textbf{3} trial $\{-1,0,+1\}$
 & \emph{degenera}: sem cúspide \hfill\textsc{expl.}
 & cúspide $\tau=0$; $+0=-0$ \hfill\textsc{expl.}
 & $\operatorname{sgn}(a/b)$ \hfill\textsc{expl.}
 & $\tau=0$ é \emph{vazio} $\iff$ irracional \hfill\textsc{expl.} \\
\textbf{4} tetral $T+T^{*}$
 & $F_{2w}=F_w\oplus\sigma F_w$ \hfill\textsc{expl.}
 & $\mathbb{N}^{2}/{\sim}$; classe $[(a,b)]$ \hfill\textsc{expl.}
 & recebida (\code{thm:juntas}) \hfill\textsc{receb.}
 & $(A,B)$, e quocientar aberto/fechado \hfill\textsc{expl.} \\
\textbf{5} pental $x^{2}=-1$
 & \emph{degenera}: uma raiz só \hfill\textsc{expl.}
 & $G_{\mathrm{geom}}\bmod2$ \hfill\textsc{receb.}
 & $J(a,b)=(-b,a)$, $J^{2}=-I$ \hfill\textsc{expl.}
 & \emph{não existe}: ordenado \hfill\textsc{expl.} \\
\textbf{6} hexal (interface)
 & base ortonormal: coordenada $=$ ler o bit \hfill\textsc{receb.}
 & $\otimes\leftrightarrow(a-b)(c-d)$ \hfill\textsc{receb.}
 & interface \hfill\textsc{receb.}
 & densidade de $\mathbb{Q}$ em $\mathbb{R}$ \hfill\textsc{receb.} \\
\textbf{7} octonião dual
 & encaixe por prefixos \hfill\textsc{expl.}
 & $q\parallel p$; par $\neq$ classe; $\Word_8$ \hfill\textsc{expl.}
 & junta dual \hfill\textsc{receb.}
 & três realizações ligadas sem fundidas \hfill\textsc{expl.} \\
\bottomrule
\end{tabular}
\end{center}
\noindent{\footnotesize \textsc{expl.} realizada neste degrau ·
\textsc{receb.} interface herdada do centro, não derivada aqui ·
\textsc{mista}/\textsc{canón.} ver o paper do degrau.
\textbf{Três leis degeneram em $k=-1$} (1, 3, 5) --- e é isso que permite a
torre binária dobrar sem excepção. A Lei~5 falta nos dois extremos por razões
\emph{opostas}: em $k=-1$ por característica~2, em $k=2$ por ordem.
\textbf{Nenhum degrau realiza as oito}; a coluna diz onde cada uma vive.\par}

\noindent As oito leis (\code{analitico sub:oitoleis}) \emph{estruturam} este
paper --- não decoram. Cada lei é primitiva com operador e medidor; neste
degrau $k{=}0$ realizam-se no molde dual $T+T^{*}$:
duplicar $\to$ ler em paralelo ($\mathrm{Dir}/\mathrm{Cruz}$) $\to$ quocientar.

\begin{center}\small
\begin{tabular}{@{}clp{4.6cm}cl@{}}
\toprule
 & lei (centro) & realização $\mathbb{N}\!\to\!\mathbb{Z}$ & estatuto & § \\
\midrule
\textbf{0} & divisão do zero & folha; $[(1,0)]\oplus[(0,1)]\sim0$; $0\dg=\infty$ elevado & explícita & \ref{sec:peano}, \ref{sec:sinal} \\
\textbf{1} & dual ($1\dg=-1$) & $\iota$; $\nu(z)=-z$; $\mathbb{Z}=\mathbb{N}+{}$sinal & explícita & \ref{sec:sinal} \\
\textbf{2} & bidual ($T\dg=-T$) & $\mathrm{Dir}/\mathrm{Cruz}$; aranha $R^{4}=\mathrm{id}$ & mista & \ref{sec:espaco}, \ref{sec:memoria} \\
\textbf{3} & trial $\{-1,0,+1\}$ & cúspide $\tau=0$; $+0=-0$; cúspide realiza $P(0)$ & explícita & \ref{sec:peano} \\
\textbf{4} & tetral ($T+T^{*}$) & $\sim$; classe $[(a,b)]$; $\Phi$; Peano & explícita & \ref{sec:peano}, \ref{sec:espaco} \\
\textbf{5} & pental ($x^{2}=-1$) & $G_{\mathrm{geom}}\bmod2$ & recebida & \ref{sec:memoria} \\
\textbf{6} & hexal (interface) & $\otimes\leftrightarrow(a-b)(c-d)$ & recebida & \ref{sec:ops} \\
\textbf{7} & octonião dual & $q$ vs.\ $p$; par $\neq$ classe; $G$ & explícita & \ref{sec:espaco}, \ref{sec:memoria} \\
\bottomrule
\end{tabular}
\end{center}

\noindent\textbf{Hierarquia} (como \code{racionais.tex}): \emph{recebe} do
\emph{Corpo Algébrico}/\emph{Analítico}; \emph{realiza} neste degrau;
\emph{verifica} nos medidores. Leis~5 e~6 não se derivam aqui --- interface
herdada. Leis~0, 1, 3, 4, 7 têm realização interna imediata em
$\mathbb{N}^{2}/{\sim}$.

%----------------------------------------------------------------------
\subsection{Justificação das oito leis neste degrau}\label{sec:leis-realizadas}

\noindent Cada entrada segue \textbf{prova} $\mid$ \textbf{realiza} $\mid$
\textbf{verifica}.

\medskip\noindent\textbf{Lei~0} (divisão do zero).
\emph{Prova:} \code{algebrico sec:lei0}, \code{thm:inversao-dual} --- o
cancelamento $1+(-1)=0$ não se postula.
\emph{Realiza:} subtracção parcial em $\mathbb{N}$ (§N0); folha
$[(1,0)]\oplus[(0,1)]\sim[(0,0)]$ induzida por $1+0=1+0$; elevação
$\infty\oplus(-1)=0$ (§M8).
\emph{Verifica:} \code{aritmetica.c} §N0; §NT5; §M8.

\medskip\noindent\textbf{Lei~1} (dual).
\emph{Prova:} \code{algebrico cor:cadeia} ($\mathbb{Z}=\mathbb{N}+{}$sinal),
\code{thm:inversao-dual} ($a-b=a\oplus b^{\dagger}$).
\emph{Realiza:} $\iota(a,b)=(b,a)$ no espaço; $\nu([(a,b)])=-(a-b)$ no
quociente; $\nu^{2}=\mathrm{id}$, $\nu(xy)=-\nu(x)\nu(y)$.
\emph{Verifica:} §NT6, §NT6b, §NT9.

\medskip\noindent\textbf{Lei~2} (bidual).
\emph{Prova:} \code{analitico sub:oitoleis} --- $T\dg=-T$ no sector
antissimétrico; \emph{não} confundir com $\iota$ (Lei~1) nem com
$V^{**}$ (degrau seguinte, \code{racionais.tex}).
\emph{Realiza:} duas réguas em paralelo --- $\mathrm{Dir}=D(a,b)=a+b$ e, no
quociente, $\mathrm{Cruz}=\Phi$ (Def.~\ref{def:espaco}); geometricamente
$T\dg=-T$ é \emph{recebida} da camada algébrica; a realização geométrica
neste degrau é $R^{4}=\mathrm{id}$ na aranha (§AG9) --- \emph{não} deduzida
de $\mathbb{N}^{2}/{\sim}$.
\emph{Verifica:} §NT10 (Dir $\neq$ invariante; Cruz $=$ invariante); §AG9.

\medskip\noindent\textbf{Lei~3} (trial).
\emph{Prova:} \code{algebrico thm:cuspide}, \code{thm:unificacao} ---
$\tau\in\{-1,0,+1\}$; centro $\tau=0$.
\emph{Realiza:} cúspide ($\tau=0$, $+0=-0$) \emph{realiza} leitura de $P(0)$;
ordens $2,3,4$ da tríade medidas (\code{lei8\_trial.c}).
\emph{Verifica:} \code{cuspide.c} §C4--§C7; \code{lei8\_trial.c} §L1--§L3.

\medskip\noindent\textbf{Lei~4} (tetral, $T+T^{*}$).
\emph{Prova:} \code{analitico thm:tecidos} --- $T_{k+1}=T_k+T_k^{*}$;
\code{algebrico def:inducao} --- indução/meta-indução.
\emph{Realiza:} duplicação $\mathbb{N}\to\mathbb{N}^{2}$; Prop.~\ref{prop:equiv}
($\sim$ equivalência); classe $[(a,b)]$; Thm.~\ref{thm:z} ($\Phi$); Peano
(Cor.~\ref{cor:peano-soma}).
\emph{Verifica:} §NP; §M0; §NT1--§NT3.

\medskip\noindent\textbf{Lei~5} (pental, $x^{2}=-1$).
\emph{Prova:} \code{algebrico} --- interface binária no centro algébrico
(\emph{não} deste paper).
\emph{Realiza:} projeção $G_{\mathrm{geom}}\mapsto G_{\mathrm{geom}}\bmod2$
na descida geométrica.
\emph{Verifica:} \code{aranha\_g.c} §AG6--§AG7.

\medskip\noindent\textbf{Lei~6} (hexal, interface).
\emph{Prova:} \code{analitico thm:estrelaseis} --- soma$=$produto é
interface do hexal (\emph{não} demonstrada neste degrau).
\emph{Realiza:} $\otimes$ bem-definido; $\mathrm{Cruz}(\otimes)=(a-b)(c-d)$
--- compatível com a interface, pré-$\mathbb{Q}$.
\emph{Verifica:} §NT4.

\medskip\noindent\textbf{Lei~7} (octonião dual, ligar sem fundir).
\emph{Prova:} \code{algebrico def:espaco} --- par distinto do quociente;
\code{arquitetura thm:multiplicidade} --- memória por duplicidade.
\emph{Realiza:} $q\colon\mathbb{N}^{2}\to\mathbb{Z}$ vs.\
$p\colon I\to\mathbb{Z}^{2}$; $G_{\mathrm{alg}}$ vs.\ $G_{\mathrm{geom}}$
(Obs.~\ref{obs:galg}); envelope $\Word_8$ restrige fibra (Def.~\ref{def:w8}).
\emph{Verifica:} §NT1--§NT2, §NT7--§NT11; §AG1--§AG11.

\medskip\noindent\textbf{RECEBE} (\code{algebrico def:espaco}): molde dual,
torre, fibra/pagamento, tríade, cúspide, Dir/Cruz.
\textbf{ACRESCENTA}: $\mathbb{N}^{2}/{\sim}\cong\mathbb{Z}$, folha, $\Phi$,
dois circuitos $q\parallel p$, envelope $\Word_8$.

%======================================================================
\section{Os naturais --- Leis 0, 3 e 4}\label{sec:peano}

\begin{definicao}[Os naturais]\label{def:peano}
$\mathbb{N}=\{0,1,2,\ldots\}$ com sucessor $S(n)=n+1$, Peano (incluso $0$):
$S(n)\neq0$; $S$ injectivo; \textbf{indução}
($P(0)\wedge\forall n(P(n)\Rightarrow P(S(n)))\Rightarrow\forall n\,P(n)$);
soma e produto recursivos usuais.
\end{definicao}

A indução de Peano fornece a prova ascendente (Thm.~\ref{thm:peano-ind}).
A meta-indução (\code{thm:meta-inducao}) é camada de \textbf{realização/verificação}
nos medidores --- não substitui o axioma de indução.

\noindent\textbf{Lei~3} (trial): a verificação pelos dois lados passa pela
tríade Cantor/Viviani/Julia (\code{algebrico thm:unificacao}) e pela
cúspide (\code{algebrico thm:cuspide}): $\tau=\operatorname{sign}(\mathrm{disc})
\in\{-1,0,+1\}$; o centro é $\tau=0$ com $+0=-0$; Cantor/Julia
($\nu^{2}=\mathrm{id}$) fixam a álgebra e trocam as pontas; o fractal de
Farey é a órbita das cúspides (\code{cuspide.c} §C4). A prova de $P(0)$
usa só Peano (Thm.~\ref{thm:peano-ind}); a cúspide ($\tau=0$, $+0=-0$)
\emph{realiza} a leitura geométrica dessa base (\code{cuspide.c} §C7).

\noindent\textbf{Lei~0} (nome do centro algébrico: \emph{divisão do zero} --- neste
degrau \textbf{não} é divisão aritmética $1/0$; é decomposição dual da classe
nula): em $(\mathbb{N},+)$ a subtracção é parcial ---
$a+x=b$ só quando $b\ge a$ (\code{aritmetica.c} §N0). A cadeia do degrau:
\[
\boxed{
\text{subtracção parcial em }\mathbb{N}
\;\longrightarrow\;
\text{duplicação }V_{\mathrm{int}}=\mathbb{N}^{2}
\;\longrightarrow\;
\text{quociente }\mathbb{N}^{2}/{\sim}
\;\longrightarrow\;
\text{subtracção total em }\mathbb{Z}
}
\]
Paga-se em $\mathbb{Z}$ pela folha (\S\ref{sec:sinal}), onde a construção
\emph{demonstra}
\[
[(1,0)]+[(0,1)]=[(1,1)]=[(0,0)]:
\]
ou, ao nível dos representantes, $(1,1)\sim(0,0)$.
Não se postula o cancelamento: é \textbf{induzido pela equivalência}
($1+0=1+0$); geometricamente, a folha $\iota$ realiza-o
(\code{algebrico thm:inversao-dual}).

\begin{teorema}[Esquema de indução]\label{thm:peano-ind}
Se $P$ verifica $P(0)$ e $\forall n\,(P(n)\Rightarrow P(S(n)))$, então
$\forall n\in\mathbb{N}\,P(n)$.
\end{teorema}

\begin{proof}
Axioma de indução de Peano (Def.~\ref{def:peano}). $\square$
\end{proof}

\begin{corolario}[Soma recursiva]\label{cor:peano-soma}
Para todo $n,m\in\mathbb{N}$: $\mathrm{nt\_soma}(n,m)=n+m$.
\end{corolario}

\begin{proof}
Indução em $m$. Base $m=0$: $\mathrm{nt\_soma}(n,0)=n=n+0$. Passo:
$\mathrm{nt\_soma}(n,S(m))=S(\mathrm{nt\_soma}(n,m))=S(n+m)=n+S(m)$ pela
definição recursiva e hipótese de indução. \emph{Prova:} só Peano
(Def.~\ref{def:peano}, Thm.~\ref{thm:peano-ind}). \emph{Realiza:} a cúspide
($\tau=0$, $+0=-0$) e a Lei~4 leem o passo dualmente
(\code{cuspide.c} §C7). \emph{Verifica:} \code{aritmetica.c} §NP;
\code{conservacao\_metrica.c} §M0. $\square$
\end{proof}

\colunas{\code{def:inducao}; \code{thm:unificacao}; \code{thm:cuspide}}
        {Peano; trial; subtracção parcial; indução em $m$}
        {§NP; §M0; \code{cuspide.c} §C4--§C7; \code{lei8\_trial.c}}

%======================================================================
\section{O espaço de pares e a classe --- Leis 2, 4 e 7}\label{sec:espaco}

\begin{definicao}[O espaço de pares e a relação]\label{def:espaco}
\textbf{Lei~4} ($T+T^{*}$): o espaço total é
\[
\boxed{\;V_{\mathrm{int}}=\mathbb{N}^{2}\;}
\]
--- duplicação de $\mathbb{N}$ --- com involução $\iota(a,b)=(b,a)$
(\textbf{Lei~1} no espaço). A relação de equivalência vive em $V_{\mathrm{int}}$:
\[
\boxed{\;(a,b)\sim(c,d)\iff a+d=b+c\;}
\]
A condição é a \textbf{soma cruzada nula}: dois naturais por lado, nenhuma
subtracção em $\mathbb{N}$.

\medskip\noindent
O inteiro é a \textbf{classe de equivalência} (como em
\code{racionais.tex} Def.~\ref{def:espaco}):
\[
\boxed{\;
\mathbb{Z}=V_{\mathrm{int}}/{\sim}
\;=\;\mathbb{N}^{2}/{\sim}\;}
\]
Elemento típico: $[(a,b)]$ --- par $\neq$ classe (\textbf{Lei~7}).

\medskip\noindent
\textbf{Duas leituras em paralelo} --- \textbf{Lei~2} e \textbf{Lei~4}
(\code{algebrico def:espaco}). \emph{Antes} do quociente não se usa
subtracção em $\mathbb{N}$:
\[
D(a,b):=a+b\qquad\text{(Dir)},\qquad
C_{\times}(a,b;c,d):\ a+d=b+c\qquad\text{(soma cruzada)}.
\]
Depois de $\Phi$ (Thm.~\ref{thm:z}), no quociente:
\[
\boxed{\;\mathrm{Cruz}(x):=\Phi(x),\qquad
D([(a,b)])=a+b\ \text{(não invariante)}.\;}
\]
Ex.: $(2,0)\sim(3,1)$ --- $\mathrm{Cruz}$ vale $2$ em ambos, mas
$D(2,0)=2\neq D(3,1)=4$ (\code{tests/inteiros.c} §NT10).
\end{definicao}

\begin{proposicao}[$\sim$ é relação de equivalência]\label{prop:equiv}
Em $V_{\mathrm{int}}=\mathbb{N}^{2}$:
\begin{enumerate}\itemsep2pt
\item \textbf{Reflexiva:} $(a,b)\sim(a,b)$ pois $a+b=b+a$.
\item \textbf{Simétrica:} $a+d=b+c\Rightarrow c+b=a+d$, logo $(c,d)\sim(a,b)$.
\item \textbf{Transitiva:} se $a+d=b+c$ e $c+f=d+e$, somando
      $a+d+c+f=b+c+d+e$; por comutatividade
      $(a+f)+(c+d)=(b+e)+(c+d)$; pela cancelabilidade em $\mathbb{N}$,
      $a+f=b+e$, logo $(a,b)\sim(e,f)$.
\end{enumerate}
\end{proposicao}

\begin{proof}
(1)--(2) imediatos. (3) Dadas $a+d=b+c$ e $c+f=d+e$, soma-se:
$a+d+c+f=b+c+d+e$. Por comutatividade,
$(a+f)+(c+d)=(b+e)+(c+d)$; pela cancelabilidade da soma em $\mathbb{N}$,
$a+f=b+e$, logo $(a,b)\sim(e,f)$. $\square$
\end{proof}

\begin{proposicao}[Soma cruzada]\label{prop:somacruz}
$(a,b)\sim(c,d)\iff C_{\times}(a,b;c,d)\iff a+d=b+c$.
\end{proposicao}

\begin{proof}
Definição de $\sim$. $\square$
\end{proof}

\begin{teorema}[Construção dos inteiros]\label{thm:z}
Seja $V_{\mathrm{int}}=\mathbb{N}^{2}$ com $(a,b)\sim(c,d)\iff a+d=b+c$
(Prop.~\ref{prop:equiv}). Então o quociente $\mathbb{Z}=V_{\mathrm{int}}/{\sim}$ recebe estrutura
de anel; o mapa
\[
\boxed{\;
\Phi\colon V_{\mathrm{int}}/{\sim}\;\overset{\cong}{\longrightarrow}\;\mathbb{Z},\qquad
\Phi([(a,b)])=a-b\;}
\]
é isomorfismo de anéis (identificação com $\mathbb{Z}$ usual). Sob $\Phi$,
\[
(a,b)\oplus(c,d)=(a+c,b+d),\qquad
(a,b)\otimes(c,d)=(ac+bd,ad+bc)
\]
(Def.~\ref{def:ops}) induzem $+$ e $\cdot$ usuais. Elementos especiais:
\[
0=[(0,0)]=[(n,n)],\qquad +1=[(1,0)],\qquad -1=[(0,1)].
\]
Além disso, $[(a,b)]=[(c,d)]\iff a-b=c-d$ e
\[
[(a,b)]\le[(c,d)]\iff a-b\le c-d
\]
(\emph{ordem transportada} por $\Phi$ --- não há ordem
componente-a-componente canônica em $\mathbb{N}^{2}$ cuja descida ao
quociente produza a ordem usual de $\mathbb{Z}$).
\end{teorema}

\begin{proof}[Prova]
(1) $\Phi$ bem-definida: Prop.~\ref{prop:somacruz}.
(2) Sobrejetiva: $k\in\mathbb{Z}$ --- se $k\ge0$, $\Phi([(k,0)])=k$; se
$k<0$, $\Phi([(0,-k)])=k$.
(3) Injectiva: $\Phi([(a,b)])=\Phi([(c,d)])$ $\Rightarrow$ $a-b=c-d$
$\Rightarrow$ $a+d=b+c$ $\Rightarrow$ $[(a,b)]=[(c,d)]$.
(4) $\Phi$ preserva $\oplus$ e $\otimes$: Thm.~\ref{thm:bem-definido} --- as
operações descem ao quociente.
(5) \textbf{Anel no quociente.} Pelo Thm.~\ref{thm:bem-definido}, $\oplus$ e
$\otimes$ são bem-definidas em $V_{\mathrm{int}}/{\sim}$ (independentes do
representante). Verificam-se os axiomas directamente no quociente:
\begin{itemize}\itemsep1pt
\item \emph{Adição:} associatividade e comutatividade herdam-se de
      $\mathbb{N}^{2}$; $[(0,0)]$ é neutro pois
      $(a,b)\oplus(0,0)=(a,b)$; oposto $[(b,a)]$ pois
      $(a,b)\oplus(b,a)=(a+b,b+a)\sim(0,0)$.
\item \emph{Multiplicação:} $\Phi([(a,b)]\otimes[(c,d)])=(a-b)(c-d)$; a
      associatividade e comutatividade de $\cdot$ em $\mathbb{Z}$ transportam-se
      por $\Phi$; $[(1,0)]$ é neutro pois $\Phi([(1,0)])=1$.
\item \emph{Distributividade:}
      $\Phi(x\otimes(y\oplus z))=\Phi(x)\bigl(\Phi(y)+\Phi(z)\bigr)
      =\Phi(x)\Phi(y)+\Phi(x)\Phi(z)
      =\Phi(x\otimes y)+\Phi(x\otimes z)$; pela injectividade de $\Phi$,
      $x\otimes(y\oplus z)=x\otimes y\oplus x\otimes z$.
\end{itemize}
Logo $(V_{\mathrm{int}}/{\sim},\oplus,\otimes)$ \textbf{é} anel. Como
$\Phi$ é isomorfismo de anéis (bijecção que preserva $+$ e $\cdot$), identifica-se
com $\mathbb{Z}$ usual.
(6) Ordem: $x\le y\iff\Phi(x)\le\Phi(y)$ com a ordem usual em $\mathbb{Z}$.
$\square$
\end{proof}

\noindent O objecto novo deste degrau não é uma nova operação sobre números: é
a representação que torna a subtracção total no quociente:
\[
\boxed{
\begin{array}{c}
\text{pares de naturais }V_{\mathrm{int}}=\mathbb{N}^{2}\\
\downarrow\\
\text{soma cruzada }a+d=b+c\\
\downarrow\\
\text{classe de equivalência }[(a,b)]\\
\downarrow\\
\mathbb{Z}=V_{\mathrm{int}}/{\sim}\\
\downarrow\\
\bar\iota=\nu\ \text{(swap)}\\
\downarrow\\
x\mapsto -x
\end{array}}
\]

\begin{corolario}[Ordem no quociente]\label{cor:ordem}
A ordem total de $\mathbb{Z}=V_{\mathrm{int}}/{\sim}$ é
\[
x\le y\iff\Phi(x)\le\Phi(y)\iff\mathrm{Cruz}(x)\le\mathrm{Cruz}(y).
\]
Ela \textbf{não} coincide com ordem componente-a-componente canônica em
$\mathbb{N}^{2}$: a ordem do quociente é a transportada por $\Phi$.
\end{corolario}

\begin{observacao}[Molde do degrau seguinte]
$\mathbb{N}\!\to\!\mathbb{Z}$: $a+d=b+c$ (soma cruzada). $\mathbb{Z}\!\to\!\mathbb{Q}$:
$ad=bc$ (produto cruzado; \code{racionais thm:q}). Analogia:
\[
\boxed{\text{soma cruzada}\longrightarrow\text{produto cruzado}}.
\]
\end{observacao}

\medskip\noindent\textbf{Lei~7} --- projeção quociente e fibra:
\[
q\colon V_{\mathrm{int}}\to\mathbb{Z},\qquad q(a,b)=[(a,b)],\qquad
q^{-1}(z)=\{(a,b)\in\mathbb{N}^{2}:a-b=z\}.
\]
Ex.: $(3,1)\sim(2,0)$ --- ambos em $q^{-1}(2)$; o par guarda
\textbf{representação}, a classe $[(a,b)]$ guarda o \textbf{valor}.
$G_{\mathrm{alg}}(z):=|q^{-1}(z)|$ (\S\ref{sec:memoria}).

\medskip\noindent\textbf{Lei~1} --- o swap desce ao quociente:
\[
\boxed{\;
\iota:V_{\mathrm{int}}\to V_{\mathrm{int}},\quad
\bar\iota:\mathbb{Z}\to\mathbb{Z},\quad
\bar\iota([(a,b)])=[(b,a)],\quad
\nu=\bar\iota,\quad
\Phi(\nu(x))=-\Phi(x).\;}
\]
Se $(a,b)\sim(c,d)$, então $a+d=b+c$ implica $b+c=a+d$, logo
$\iota(a,b)\sim\iota(c,d)$: $\bar\iota$ é bem-definida
(\code{racionais.tex} §\ref{sec:espaco}, mesmo argumento).

\colunas{\code{algebrico def:espaco}; \code{racionais def:espaco}}
        {$\sim$ equivalência; classe $[(a,b)]$; $\Phi=a-b$; $\bar\iota$}
        {§NT1--§NT2, §NT10}

%======================================================================
\section{Operações e folha --- Leis 1 e 6}\label{sec:ops}

\begin{definicao}[Operações no par]\label{def:ops}
Operações em $V_{\mathrm{int}}=\mathbb{N}^{2}$ \emph{antes} do quociente:
\[
(a,b)\oplus(c,d):=(a+c,b+d),\qquad
(a,b)\otimes(c,d):=(ac+bd,ad+bc),\qquad
\boxed{\;\ominus(a,b):=\iota(a,b)=(b,a).\;}
\]
$\mathbb{N}^{2}$ com $\oplus$ e $\otimes$ \textbf{não} é anel --- a
estrutura desce a $V_{\mathrm{int}}/{\sim}\cong\mathbb{Z}$ (Thm.~\ref{thm:z}).
\end{definicao}

\noindent\textbf{Lei~6} (interface recebida): no quociente, com
$\mathrm{Cruz}=\Phi$,
\[
\boxed{\;
\mathrm{Cruz}(x\cdot y)=\mathrm{Cruz}(x)\cdot\mathrm{Cruz}(y)\;}
\]
--- pois $\Phi(xy)=\Phi(x)\Phi(y)$; a identidade ``soma$=$produto'' pertence
ao degrau anterior.

\begin{teorema}[Bem-definido]\label{thm:bem-definido}
Se $(a,b)\sim(a',b')$, $(c,d)\sim(c',d')$, então $\oplus$ e $\otimes$
descendem a $+$ e $\cdot$ em $\mathbb{Z}$; $\iota(a,b)=(a,b)\iff a=b$.
\end{teorema}

\begin{proof}
\textbf{(1) Soma:} se $(a,b)\sim(a',b')$ e $(c,d)\sim(c',d')$, então
$a+b'=b+a'$ e $c+d'=d+c'$. Somando:
\[
(a+c)+(b'+d')=(a+b')+(c+d')=(b+a')+(d+c')=(a'+c')+(b+d),
\]
logo $(a+c,b+d)\sim(a'+c',b'+d')$.
\textbf{(2) Produto:} se $(a,b)\sim(a',b')$ e $(c,d)\sim(c',d')$, então
$a-b=a'-b'$ e $c-d=c'-d'$. Logo
\[
(ac+bd)-(ad+bc)=(a-b)(c-d)=(a'-b')(c'-d')=(a'c'+b'd')-(a'd'+b'c'),
\]
donde $(ac+bd)+(a'd'+b'c')=(a'c'+b'd')+(ad+bc)$, ou seja,
$(ac+bd,ad+bc)\sim(a'c'+b'd',a'd'+b'c')$. Também
$\Phi([(a,b)]\otimes[(c,d)])=(a-b)(c-d)=\Phi(x)\Phi(y)$ para
$x=[(a,b)]$, $y=[(c,d)]$.
\textbf{(3)} $\iota(a,b)=(a,b)\iff a=b$. $\square$
\end{proof}

\begin{proposicao}[Fixos]\label{prop:fixos}
$\operatorname{Fix}(\iota)=\{(n,n)\}$ no espaço;
$\operatorname{Fix}(\nu)=\{0_{\mathbb{Z}}\}$ no quociente.
\end{proposicao}

\begin{corolario}[Embedding e folha]\label{cor:folha}
$\eta(n)=[(n,0)]$ injectivo; $0_{\mathbb{Z}}=[(n,n)]$;
oposto $=[(b,a)]$. Folha: $(n,0)\mapsto(0,n)$ e
$[(n,0)]\oplus[(0,n)]\sim[(0,0)]$.
\end{corolario}

\colunas{\code{thm:operacao}}
        {$\oplus$, $\otimes$, $\ominus$}
        {§NT3--§NT5}

%======================================================================
\section{Sinal e involução --- Leis 0 e 1}\label{sec:sinal}

\noindent\textbf{Lei~1}: $\mathbb{Z}=\mathbb{N}+{}$sinal
(\code{algebrico cor:cadeia}). O sinal é escolha de metade após o espelho:
$+1=[(1,0)]$, $-1=[(0,1)]$, $a-b=a\oplus b^{\dagger}$ com $1^{\dagger}=-1$
(\code{algebrico thm:inversao-dual}).

\begin{definicao}[Involuções --- três papéis distintos]\label{def:nu}
\[
\boxed{
\iota\ \text{(Lei~1, espaço)},\quad
\bar\iota=\nu\ \text{(Lei~1, quociente)},\quad
\dagger\ \text{(Lei~0/1, operação recebida)}.
}
\]
$\iota(a,b)=(b,a)=\ominus(a,b)$; $\nu([(a,b)])=\bar\iota([(a,b)])=[(b,a)]$;
$\nu=\bar\iota=q\circ\iota$ (§\ref{sec:espaco}). Dual no espaço
$\xrightarrow{q}$ inverso aditivo no quociente:
\[
\boxed{
\begin{array}{ccc}
V_{\mathrm{int}} & \xrightarrow{\;\iota\;} & V_{\mathrm{int}} \\
\downarrow q && \downarrow q \\
\mathbb{Z} & \xrightarrow{\;\nu\;} & \mathbb{Z}
\end{array}
\qquad q\circ\iota=\nu\circ q.}
\]
\end{definicao}

\begin{teorema}[$\nu$ --- Lei~1]\label{thm:nu}
Com $\nu=-\mathrm{id}_{\mathbb{Z}}$: $\nu\circ\nu=\mathrm{id}$;
$\mathrm{Fix}(\nu)=\{0\}$; $\nu(x+y)=\nu(x)+\nu(y)$;
$\boxed{\nu(xy)=-\nu(x)\nu(y)}$.
Portanto $\nu$ é automorfismo \emph{aditivo} involutivo --- \emph{não} é,
em geral, automorfismo multiplicativo de anel.
\end{teorema}

\begin{proof}
(1)--(2) directos do swap. (3) como antes. (4)
$\nu(xy)=-(xy)$ e $-\nu(x)\nu(y)=-(-x)(-y)=-(xy)$. $\square$
\end{proof}

\begin{teorema}[Mediante --- elevação de $\infty+1=-1$]\label{thm:mediante}
$1+(-1)=0$. A componente $0\dg=\infty$ pertence à \textbf{interface recebida}
da camada projectiva/Möbius (\code{algebrico sec:lei0}) --- \emph{não} é
necessária para $\mathbb{N}^{2}/{\sim}\cong\mathbb{Z}$. No sector
projectivo, a interface possui $M(\infty)=-1$ e $\infty\oplus(-1)=0$
(\code{descida\_mobius.c} §M8). \textbf{Neste degrau}, a realização
independente da mesma folha é
$[(1,0)]+[(0,1)]=[(0,0)]$, sem dependência lógica do sector projectivo.
\end{teorema}

\colunas{\code{cor:cadeia}; \code{sec:lei0}}
        {$\nu$; folha; $+1,-1,0$}
        {§NT5--§NT6b, §NT9; §M8}

%======================================================================
\section{Memória por duplicidade --- Leis 2, 5 e 7}\label{sec:memoria}

\noindent\textbf{Lei~7} materializa-se dinamicamente
(\code{arquitetura thm:multiplicidade}). Realização geométrica
$p\colon I\to\mathbb{Z}^{2}$, $G_{\mathrm{geom}}(x):=|p^{-1}(x)|$:
\[
\boxed{
\text{trajetória}\xrightarrow{\,p\,}\text{dobra}\xrightarrow{|p^{-1}|}
G_{\mathrm{geom}}\xrightarrow{\text{leitura}}\text{percepção}.
}
\]
Duplicidade: $p(i)=p(j)$ com $i\neq j$ $\Rightarrow$ $G_{\mathrm{geom}}(p(i))>1$.
Memória no \textbf{espaço}: $G_{t+1}(x)=G_{t}(x)+\mathbf{1}_{\{x=p(t)\}}$.

\medskip\noindent\textbf{Realização com dobra.} O teorema é sobre $p$ e
$G_{\mathrm{geom}}$; a curva do dragão (Heighway) é \textbf{teste} recursivo
de dobra (\code{aranha\_g.c} §AG2). Recta injectiva: $G_{\mathrm{geom}}\equiv1$
(§AG3) --- controlo sem identificação de índices.

\medskip\noindent\textbf{Aranha estigmérgica}: lê $G_{\mathrm{geom}}$ (§AG8),
escreve sem memória interna (§AG10). \textbf{Lei~2}:
$T\dg=-T\leadsto R^{4}=\mathrm{id}$ --- realização geométrica, não
identidade formal (§AG9).

\medskip\noindent\textbf{Lei~5} (interface recebida do \emph{Corpo Algébrico}):
$x^{2}=-1$ fornece a interface binária no centro algébrico; aqui
realiza-se a projeção $G_{\mathrm{geom}}\mapsto G_{\mathrm{geom}}\bmod 2$
(§AG6--§AG7) --- \emph{não} é derivação deste paper.

\medskip\noindent\textbf{Dois padrões, dois espaços.}
\begin{enumerate}\itemsep2pt
\item \textbf{Fibra algébrica}: $q(a,b)=[(a,b)]$,
      $q^{-1}(z)=\{(a,b):a-b=z\}$, $G_{\mathrm{alg}}(z)=|q^{-1}(z)|$.
      Ex.: $(3,1)\sim(2,0)$ --- o par guarda \textbf{representação}; a
      classe guarda o \textbf{valor quocientado}.
\item \textbf{Fibra geométrica}: $p(i)=x$,
      $p^{-1}(x)=\{i\in I:p(i)=x\}$, $G_{\mathrm{geom}}(x)=|p^{-1}(x)|$.
      O índice/célula fornece o \textbf{endereço} espacial.
\end{enumerate}
\[
\boxed{
G_{\mathrm{alg}}\neq G_{\mathrm{geom}},\qquad
G_{\mathrm{alg}}(z)=\aleph_0\ \text{(fibra estrutural)},\qquad
G_{\mathrm{geom}}(x)=|p^{-1}(x)|\ \text{(multiplicidade da trajetória)}.
}
\]
\begin{observacao}[$G_{\mathrm{alg}}$ no espaço ilimitado]\label{obs:galg}
Em $\mathbb{N}^{2}$ sem envelope, $q^{-1}(z)=\{(a,b):a-b=z\}$ é
\textbf{infinita enumerável} para todo $z$ (ex.: $z=2$:
$(2,0),(3,1),(4,2),\ldots$). Logo $G_{\mathrm{alg}}(z)=\aleph_0$:
cardinalidade \emph{estrutural}, não contador finito de memória --- a
fibra algébrica não é «memória acumulada»; é propriedade cardinal de $q$.
A memória operacional é $G_{\mathrm{geom}}$ (multiplicidade temporal de $p$,
ou fibra restrita a $\Word_8$).
\end{observacao}
Mesmo padrão estrutural (duplicar $\to$ identificar $\to$ contar $\to$
projectar); espaços distintos.

\subsection{Par homomorfismo/unimodular e conservação de área}\label{sec:metrica-dragao}

O \emph{Corpo Algébrico} (\code{algebrico thm:metrica}) distingue três
camadas: a métrica canónica de Lebesgue
$d(A,B)=\mu(A\triangle B)$ --- no espaço \textbf{discreto}, preservar a
métrica reduz-se a preservar a medida de contagem: uma bijecção do conjunto
finito é \textbf{isometria discreta} ($d(TA,TB)=d(A,B)$), e \emph{não} se
confunde com homeomorfismo contínuo (no finito toda permutação é homeomorfismo
e não distingue métrica nenhuma; cf.\ \code{algebrico thm:metrica}).
Já a \textbf{conservação da área} exige o jacobiano: $\mu(TA)=|\det T|\,\mu(A)$.
No reticulado, $[\mathbb{Z}^{2}:T(\mathbb{Z}^{2})]=|\det T|$ (\code{metrica.c} §M8).
O \emph{Corpo Topológico} (\code{topologico thm:det-volume}) identifica
$|\det A_m|=1$ com o factor de volume do gato.

\medskip\noindent\textbf{Neste degrau}, o par algébrico/topológico
realiza-se em $\mathbb{Z}^{2}$:
\[
\boxed{
\begin{array}{ll}
\textbf{HOMOMORFISMO} & \Phi:V_{\mathrm{int}}/{\sim}\overset{\cong}{\to}\mathbb{Z}\ \text{(anel)};\\
& \nu=-\mathrm{id}_{\mathbb{Z}}\ \text{--- automorfismo aditivo, não de anel}\\[4pt]
\textbf{ISOMETRIA DISCRETA} & d(A,B)=\mu(A\triangle B)\ \text{no toro finito:}\\
& \text{bijecção preserva contagem}\Rightarrow\text{isometria (§M1--§M2)}\\[4pt]
\textbf{UNIMODULAR} & T\in\mathrm{GL}(2,\mathbb{Z}),\ |\det T|=1:\ \text{automorfismo}\\
& \text{do reticulado }\mathbb{Z}^{2};\ \text{a extensão }\mathbb{R}\text{-linear}\\
& T_{\mathbb{R}}\colon\mathbb{R}^{2}\to\mathbb{R}^{2}\ \text{é difeomorfismo}\\
& \text{(jacobiano constante }\det T=\pm1\text{) que preserva área;}\\
& \mu(TA)=\mu(A)\ \text{(§M3, §M8)}
\end{array}}
\]

\medskip\noindent\textbf{Dragão como teste de dobra.} A projeção
$p\colon I\to\mathbb{Z}^{2}$ percorre o reticulado; cada passo é
translação ($|\det|=1$). A curva do dragão (Heighway) \textbf{dobra} e
produz $G_{\mathrm{geom}}(x)>1$ (§AG2): duplicidade de pré-imagem no
espaço, sem perder contagem total ($\sum_x G_{\mathrm{geom}}(x)=|I|$, §AG5).
A conservação de \textbf{área discreta} (células) sob mapas unimodulares
é orthogonal à fibra algébrica $G_{\mathrm{alg}}=\aleph_0$: o dragão
testa \textbf{memória geométrica por dobra}, não o cardinal de
$q^{-1}(z)$. A descida $G_{\mathrm{geom}}\bmod 2$ (§AG6--§AG7) perde
profundidade --- exactamente como a interface binária do centro algébrico.

\colunas{\code{algebrico thm:metrica}; \code{topologico thm:det-volume}}
        {$\Phi$; isometria $d=\mu(\triangle)$; unimodular $|\det|=1$; dragão}
        {\code{metrica.c} §M3--§M8; \code{aranha\_g.c} §AG2, §AG5--§AG7}

\colunas{\code{thm:multiplicidade}; \code{def:peano}}
        {$G_{\mathrm{geom}}=|p^{-1}|$; dragão como teste}
        {\code{aranha\_g.c} §AG1--§AG11}

%======================================================================
\section{Envelope $\Word_8$ --- Lei~7 na descida}\label{sec:palavra8}

\noindent\textbf{Lei~7} restringe a fibra a um envelope finito; \textbf{Lei~4}
materializa $T+T^{*}$ em representação computável antes do overflow.

\noindent $\Word_8=\{0,\ldots,255\}$: \textbf{envelope finito de 8 bits}
dos naturais --- $\mathbb{N}$ não se redefine como $\Word_8$.
$\operatorname{succ}_8(w):=(w+1)\bmod 256$; wrap contado (\code{lib/naturais.h},
\code{nt\_wrap}) --- \emph{não} é Peano. Cadeia:
$\Word_8\to E_{16}\to\cdots\to\mathrm{GF}(2)$
(\code{arquitetura sec:descida}).

\begin{definicao}[Leitura cruzada exacta e projeções]\label{def:w8}
Par $(a,b)\in\Word_8^{2}$. A equivalência $(a,b)\sim(c,d)$ testa-se pela
\textbf{leitura cruzada exacta} em $\mathbb{N}$:
\[
\widetilde{C}(a,d):=a+d\in\{0,\ldots,510\}
\quad\text{(uint16, como \code{d16\_mult} em \code{racionais.tex})}.
\]
$(a,b)\sim(c,d)\iff \widetilde{C}(a,d)=\widetilde{C}(b,c)$ --- \emph{sem}
projectar para o byte. Quando o resultado sai do byte, o hot path (20/08) é
\textsc{saturo}$\to$\textsc{promove}:
\[
\pi_{\mathrm{wrap}}(s):=s\bmod 256,\qquad
\pi_{\mathrm{prom}}(s):=s\in\{0,\ldots,65535\}
\quad\text{(\code{uint16} exacto)}.
\]
A antiga $\pi_{\mathrm{sat}}(s)=\min(s,255)$ ficou só como vista legado em
\code{W8X.estreito}. Contadores \code{w8\_wrap} e \code{w8\_saturou} registam
saídas do envelope \textbf{à parte dos defeitos} (cf.\ \code{qz\_saturou}).
\end{definicao}

\begin{proposicao}[Políticas de overflow]\label{prop:perde8}
$\pi_{\mathrm{wrap}}$ e $\pi_{\mathrm{prom}}$ são construções distintas:
para $s>255$, $\pi_{\mathrm{wrap}}(s)\neq\pi_{\mathrm{prom}}(s)$ em geral
(ex.: $s=300$ dá $44$ vs.\ $300$). A equivalência $\sim$ em $\Word_8^{2}$
\textbf{não} depende de nenhuma das duas --- usa $\widetilde{C}$ em uint16
(§NT8). Só a \emph{escrita} no envelope escolhe política; o hot path promove.
\end{proposicao}

\medskip\noindent\textbf{Implementação} (\code{lib/naturais.h}):
\begin{center}\small
\begin{tabular}{@{}llll@{}}
\toprule
operação & função & o que faz & nota \\
\midrule
equiv.\ cruzada & \code{w8\_equiv} & $a+d=b+c$ em uint16 & não projecta \\
projeção wrap & \code{w8\_proj\_wrap} & $s\bmod 256$ & conta \code{w8\_wrap} \\
projeção/promove & \code{w8\_proj\_sat} & exacto em uint16 & conta \code{w8\_saturou} \\
migração vista & \code{W8X}/\code{w8\_x*} & estreito$=$sat legado & \code{.promovido} exacto \\
sucessor byte & \code{nt\_S} em $\Word_8$ & wrap $255\to0$ & conta \code{nt\_wrap} \\
\bottomrule
\end{tabular}
\end{center}

\noindent Hot path (20/08): \textsc{saturo}$\to$\textsc{promove} ---
\code{w8\_proj\_sat} já não devolve $\min(s,255)$; \code{qz()} já não devolve
$\pm32767$ (\code{arquitetura sec:pipelines}; \code{tests/promocao.c}).

\medido{\code{inteiros.c} §NT1--§NT11; \code{aranha\_g.c} §AG6--§AG7;
\code{cadeia.c} §KW ($10/10$); \code{migracao.c} §GW ($7/7$);
\code{reducao.c} §DW; \code{dual16.c} §EW; \code{palavra.c} §PW;
\code{palavra8.c} §P0--§P4 ($5/5$): $\oplus$/$\wedge$, $F_8$, subida a $\mathbb{N}$,
largura $8\to16\to32$, $\sim$ em $\Word_8^{2}$ via \code{lib/palavra8.h}.}

%======================================================================
\section{Fundamentação --- inventário das oito leis}\label{sec:fundamentacao}

\begin{center}\scriptsize
\begin{tabular}{@{}cllll@{}}
\toprule
lei & enunciado (centro) & realização & medidor & \\
\midrule
0 & $0=(+1)\oplus(-1)$; $0\dg=\infty$
  & subtracção parcial; folha; §M8 & §N0; §NT5; §M8 & $\checkmark$ \\
1 & $1\dg=-1$; $\nu^{2}=\mathrm{id}$
  & $\iota$; $\nu$; $\mathbb{N}+{}$sinal & §NT6; §NT6b; §NT9 & $\checkmark$ \\
2 & $T\dg=-T$
  & Dir/Cruz; $R^{4}=\mathrm{id}$ aranha & §NT10; §AG9 & $\checkmark$ \\
3 & trial $\{-1,0,+1\}$
  & cúspide; $+0=-0$; ordens 2,3,4 & §C4--§C7; §L1--§L3 & $\checkmark$ \\
4 & $T+T^{*}$
  & $\sim$ equivalência; classe; Thm.\ $\Phi$; Peano & §NP; §M0; §NT1--§NT3 & $\checkmark$ \\
5 & $x^{2}=-1$ (recebida)
  & $G_{\mathrm{geom}}\bmod2$ & §AG6--§AG7 & $\checkmark$ \\
6 & hexal (recebida)
  & $\otimes$; $(a-b)(c-d)$ & §NT4 & $\checkmark$ \\
7 & ligar sem fundir
  & $q\parallel p$; $G$; $\Word_8$ & §NT1--§NT2; §NT7--§NT11; §AG1--§AG11 & $\checkmark$ \\
\bottomrule
\end{tabular}
\end{center}

\noindent Nenhuma linha substitui demonstração matemática: verifica realização
finita nos domínios indicados (\code{racionais.tex}, §Resultados já obtidos).
Leis~5--6: interface \emph{recebida}, não derivada aqui.

%======================================================================
\section*{Síntese --- oito leis, dois circuitos}

\begin{observacao}[o Corpo Inteiro, e o que não é]
\texttt{Corpo Inteiro} é o nome estrutural desta obra. \textbf{Estatuto:}
domínio de integridade (anel no quociente; no alvo de $\Phi$, $\mathbb{Z}$
usual). Não é corpo. A passagem $\mathbb{N}\to\mathbb{Z}$ fecha por
\textbf{opostos}, não por $\mathcal{D}$.

$\eta(n)=[(n,0)]$ é inclusão, não isomorfismo nem dualidade.
$\nu(z)=-z$ é involução de $(\mathbb{Z},+)$, $\nu\neq\mathcal{D}$, $\nu\neq\star$.
$\operatorname{Star}(\mathbb{Z})=\varnothing$. Completude: \texttt{N/A}.
A passagem a $\mathbb{Q}$ é localização (quocientes), não dual.
\end{observacao}

\noindent Cada lei com prova $\mid$ realiza $\mid$ verifica
(§\ref{sec:leis-realizadas}):

\begin{itemize}\itemsep2pt
\item[\textbf{0}] Fibra parcial em $\mathbb{N}$; folha induz $0$; $\infty\oplus(-1)=0$.
\item[\textbf{1}] $\iota$ no par; $\nu(z)=-z$; $\nu(xy)=-\nu(x)\nu(y)$.
\item[\textbf{2}] $T\dg=-T$ recebida; realização $R^{4}=\mathrm{id}$ na aranha (§AG9).
\item[\textbf{3}] Trial $\tau=0$; cúspide \emph{realiza} leitura de $P(0)$; ordens 2,3,4.
\item[\textbf{4}] $T+T^{*}$: $\sim$; classe $[(a,b)]$; $\Phi=a-b$; Peano.
\item[\textbf{5}] Interface recebida: $G_{\mathrm{geom}}\bmod2$.
\item[\textbf{6}] Interface recebida: $\otimes\leftrightarrow(a-b)(c-d)$.
\item[\textbf{7}] $q\parallel p$; representante $\neq$ classe; memória $=|p^{-1}|$; $\Word_8$.
\end{itemize}

\medskip\noindent\textbf{Dois circuitos paralelos:}
\[
\mathbb{N}\xrightarrow{\text{duplicação}}\mathbb{N}^{2}
\xrightarrow{q}\mathbb{Z}
\qquad
I\xrightarrow{p}\mathbb{Z}^{2}
\xrightarrow{|p^{-1}|}G_{\mathrm{geom}}
\xrightarrow{\bmod 2}\mathbb{Z}/2\mathbb{Z}\cong\mathrm{GF}(2).
\]
\[
\boxed{
\begin{array}{c}
\mathbb{N}\xrightarrow{\text{duplicar}}\mathbb{N}^{2}
\xrightarrow{\;q\;} V_{\mathrm{int}}/{\sim}\;\cong\;\mathbb{Z}\\[2pt]
(a,b)\sim(c,d)\iff a+d=b+c,\qquad
\Phi([(a,b)])=a-b
\end{array}}
\]
\[
\boxed{
[(a,b)]\xrightarrow{\ \nu\ }[(b,a)],\qquad
\Phi(\nu(x))=-\Phi(x),\qquad
[(n,0)]\oplus[(0,n)]\sim[(0,0)].
}
\]
Descida para 8 bits: equivalência por uint16 exacto; wrap e saturação
contados à parte na escrita (Def.~\ref{def:w8}).

\medskip\noindent
\code{python3 tools/refcruz.py} (chave \code{inteiros}) ·
\code{racionais.tex} · \code{corpo\_algebrico.tex} · \code{arquitetura.tex}.

\vfill
{\footnotesize\color{cinza}\noindent Proprietário --- uso mediante contrato e pagamento (ver \texttt{LICENSE}).\par}

\end{document}
